
\end{equation}
où $w\in\Filt_{j-1}(X)$. Par l'hypothèse de récurrence on a $x\cdot w\in\Filt_{j-i-1}(X)$, d'où, compte tenu des formules (1.3.8) et (1.3.9), l'assertion
\[
x\cl_\bullet(\underline{O}_X)\in\Filt_{j-i}(X),\qquad\text{C.Q.F.D.}
\]

\textbf{Remarque 1.4.} Soit $X$ un schéma localement noethérien. On notera que nous avons fait usage des \emph{seules} propriétés suivantes de la « fonction dimension » $D(x)=D_Y(x)=\dim(\overline{\{x\}})$ sur les $X$-schémas propres $Y$ au cours de cette section :
\begin{enumerate}[label=(\roman*)]
\item Pour tout $X$-schéma propre $Y$ et pour tout $x\in Y$, on a $D_Y(x)\in\mathbb Z$.
\item Si $y\ne x$ est une spécialisation de $x$, on a
\[
D(y)<D(x).
\]
\item Pour tout $X$-morphisme $f:Y_1\to Y_2$ de $X$-schémas propres et pour tout $y\in Y_1$, on a
\[
D_{Y_2}(f(y))\leq D_{Y_1}(y).
\]
\item Il y a égalité dans (iii), si $f$ est une immersion ou si $Y_1$ est irréductible, $f$ birationnel et surjectif, et $y$ le point générique de $Y_1$.
\end{enumerate}

\uline{Cela signifie que tous les résultats de cette section, en particulier} (1.3.2) \uline{et} (1.3.3)\uline{, restent valables tels quels, si l'on y remplace la} « \uline{dimension} » \uline{par une fonction} $D$ \uline{satisfaisant aux conditions} (i)à(iv) \uline{ci-dessus et si l'on définit la filtration de} $K_\bullet(X)$ \uline{suivant une telle fonction} $D$.

